\documentclass[10pt]{article}
\usepackage[T1]{fontenc}
\usepackage[utf8]{inputenc}
\usepackage[french]{babel}
\usepackage{amsmath,amssymb,mathrsfs,mathpazo}
\usepackage{geometry}
\geometry{a4paper,margin=2cm}
\usepackage{graphicx,array,tabularx,tabularray,colortbl,multirow,multicol,booktabs,enumitem,pifont,ulem,fancybox,diagbox,cancel}
\usepackage{pstricks,pstricks-add,pst-plot,pst-tree,pst-node,pst-text,pst-func,pst-math,pst-bspline,pst-3dplot,pst-eucl,pst-solides3d}
\let\pscalculate\undefined
\usepackage{pst-bezier}
\usepackage{tikz,pgfplots,tkz-tab}
\usetikzlibrary{arrows,arrows.meta,patterns,calc,positioning,trees,decorations.markings,shapes.geometric,shapes.misc,shapes.arrows,intersections,decorations.pathmorphing,decorations.pathreplacing,backgrounds,shadows,babel}
\pgfplotsset{compat=1.18}
\usepgfplotslibrary{fillbetween}
\usepackage[np]{numprint}
\usepackage{listings,makecell,pas-tableur,twoopt}
\let\barre\undefined
\usepackage[pstricks]{tablvar}
\definecolor{darkgreen}{RGB}{0,100,0}
\providecommand{\red}{\color{red}}
\providecommand{\blue}{\color{blue}}
\providecommand{\green}{\color{green}}
\providecommand{\black}{\color{black}}
\providecommand{\white}{\color{white}}
\providecommand{\gray}{\color{gray}}
\providecommand{\cyan}{\color{cyan}}
\providecommand{\magenta}{\color{magenta}}
\providecommand{\yellow}{\color{yellow}}
\providecommand{\violet}{\color{violet}}
\providecommand{\orange}{\color{orange}}
\providecommand{\euro}{\texteuro}
\providecommand{\EUR}{\texteuro}
\providecommand{\ieme}{\textsuperscript{e}}
\providecommand{\degres}{\ensuremath{^\circ}}
\providecommand{\degree}{\ensuremath{^\circ}}
\providecommand{\vv}[1]{\overrightarrow{#1}}
\providecommand{\Coord}[1]{#1}
\providecommand{\decosix}{}
\providecommand{\starredbullet}{\ensuremath{\bullet}}
\pagestyle{empty}
\setlength{\parindent}{0pt}
\usepackage{adjustbox}\DeclareUnicodeCharacter{2236}{:}\DeclareUnicodeCharacter{2190}{\ensuremath{\leftarrow}}
\providecommand{\R}{\mathbb{R}}
\providecommand{\N}{\mathbb{N}}
\providecommand{\Z}{\mathbb{Z}}
\providecommand{\Q}{\mathbb{Q}}
\providecommand{\D}{\mathbb{D}}
\providecommand{\C}{\mathbb{C}}
\providecommand{\e}{{\mathrm{e}}}
\providecommand{\vect}[1]{\overrightarrow{#1}}
\providecommand{\vectt}[1]{\overrightarrow{\mathrm{#1}}}
\providecommand{\vvt}[1]{\overrightarrow{\mathrm{#1}}}
\providecommand{\np}[2][]{{#2~#1}}
\providecommand{\vv}[1]{{\overrightarrow{#1}}}
\renewcommand{\d}{{\mathrm{d}}}
\providecommand{\degres}{{^\circ}}
\providecommand{\degre}{{^\circ}}
\providecommand{\blue}{\color{blue}}
\providecommand{\red}{\color{red}}
\providecommand{\enskip}{\quad}
\providecommand{\llbracket}{\lbrack\!\lbrack}
\providecommand{\rrbracket}{\rbrack\!\rbrack}
\providecommand{\euro}{{\text{€}}}
\providecommand{\no}{{\text{n°}}}
\providecommand{\interval}[2]{{[#1;#2]}}
\providecommand{\intervFO}[2]{{[#1;#2[}}
\providecommand{\dotfill}{\ldots}

\providecommand{\ds}{\displaystyle}
\providecommand{\textding}[1]{#1}
\providecommand{\Oij}{\(\mathrm{O};\vec{\imath},\vec{\jmath}\)}
\providecommand{\Oijk}{\(\mathrm{O};\vec{\imath},\vec{\jmath},\vec{k}\)}
\providecommand{\Ouv}{\(\mathrm{O};\vec{u},\vec{v}\)}
\providecommand{\intervFF}[2]{\ensuremath{[#1;#2]}}
\renewcommand{\euro}{\ifmmode\text{\texteuro}\else\texteuro\fi}
\DeclareUnicodeCharacter{00B0}{\ifmmode^{\circ}\else\textdegree\fi}
\DeclareUnicodeCharacter{20AC}{\ifmmode\text{\texteuro}\else\texteuro\fi}
% Énoncés : Éducation nationale. Transcription et corrigés : APMEP, Denis Vergès et contributeurs. Découpage et classement : CoopMaths, Sébastien Lozano, Éric Elter et contributeurs. Adaptation de la présentation : Maths Explorer.
% CC BY-SA : https://coopmaths.fr/www/about/
% https://coopmaths.fr/alea/static/e3c/2021/tex/e3c_2021_06_ameriquenord_2_cor.tex
\begin{document}
\noindent\textit{APMEP / CoopMaths. CC BY-SA. Présentation adaptée par Maths Explorer.}\par\medskip
\medskip

La population d'une ville A augmente chaque année de 2\,\%. 
La ville A avait \np{4600} habitants en 2010.
La population d'une ville B augmente de 110 habitants par année.
La ville B avait \np{5100} habitants en 2010.

Pour tout entier $n$, on note $u_n$ le nombre d'habitants de la ville A et $v_n$ le nombre d'habitants de la ville B à la fin de l'année $2010 + n$.

\medskip

\begin{enumerate}
\item
\begin{list}{\textbullet}{À la fin de 2011:}
\item Le nombre d'habitants de la ville A est: \\[5pt]
$u_1=u_0 + u_0\times \dfrac{2}{100}= \np{4600} + \np{4600}\times \dfrac{2}{100}=  \np{4692}$.

\item Le nombre d'habitants de la ville B est: $v_1=v_0+110 = \np{5100}+110 = \np{5210}$.
\end{list}

\item
\begin{list}{\textbullet}{}
\item Ajouter $2\,\%$, c'est multiplier par $1+\dfrac{2}{100}=1,02$; donc la suite $(u_n)$ est géométrique de raison $q=1,02$ et de premier terme $u_0=\np{4600}$.
\item On passe de $v_n$ à $v_{n+1}$ en ajoutant $110$, donc la suite $(v_n)$ est arithmétique de raison $r=110$ et de premier terme $v_0=\np{5100}$.
\end{list}

\item 
La suite $(u_n)$ est géométrique de raison $q=1,02$ et de premier terme $u_0=\np{4600}$ donc, pour tout entier naturel $n$, on a: $u_n=u_0\times q^n = \np{4600}\times 1,02^n$.

Le nombre d'habitants dans la ville A en 2020 est: $u_{10}=\np{4600}\times 1,02^{10}\approx \np{5607}$.

\item 
La suite $(v_n)$ est arithmétique de raison $r=110$ et de premier terme $v_0=\np{5100}$ donc pour tout entier naturel $n$, on a:
$u_n=u_0 + n\times r = \np{5100} + 110n$.

Le nombre d'habitants dans la ville B en 2020 est: $v_{10}= \np{5100}+110\times 10 = \np{6200}$.

\item On complète l'algorithme ci-dessous qui permet de déterminer au bout de combien d'années la population de la ville A dépasse celle de la ville B.

\begin{center}
\renewcommand{\arraystretch}{1.2}
\fbox{
\begin{tabular}{l}
def année () : \hspace*{3cm}\\
\quad $u =\np{4600}$\\
\quad $v =\np{5100}$\\
\quad $n = 0$\\
\quad while $\blue u \leqslant v$\\
\qquad $u = \blue u \times 1,02$\\
\qquad $u = \blue v + 110$\\
\qquad $n = \blue n+1$\\
\quad return $n$\\ 
\end{tabular}
}
\end{center}
\end{enumerate}

\bigskip

\end{document}
